\section{Introduction}\label{sec:intro}

Large language models (LLMs) are now used routinely to label text at a scale and cost that
human coding cannot match, in political science, management research and public
administration among other fields \citep{Gilardi2023,Carlson2025}. Each label, however, is
not the output of ``the model'' but of a \emph{configuration}: a particular model, a
particular prompt and a particular stochastic decoding run. Semantically equivalent prompts
change accuracy by tens of percentage points \citep{Barrie2024,Chen2026}, nominally
deterministic decoding is not deterministic \citep{Atil2024}, and the choice of model,
prompt and temperature can flip the substantive conclusions of a downstream analysis
\citep{Baumann2025}. Treating LLM output as a fixed measurement therefore hides a large
and structured source of error.

The natural response is to annotate every item with several configurations and aggregate.
The dominant aggregation model, due to \citet{Dawid1979}, treats each annotator as an
independent noisy channel with its own class-conditional error rates and estimates true
labels and error rates jointly without a gold standard. Its assumption of conditional
independence given the true class was reasonable for human coders working separately. It
is not reasonable for an LLM ensemble: configurations that share a model share its
representation of the text; configurations that share a prompt share its reading of the
task; and all configurations share the training corpora and alignment procedures of the
model families that produced them. The diagnostic-testing literature learned decades ago
that ignoring such dependence overstates the accuracy of the individual tests and biases
prevalence estimates \citep{Vacek1985,Qu1996}.

This paper proposes a latent class model in which the dependence between LLM annotations
is not assumed away, nor left arbitrary, but derived from the factorial design of the
annotation ensemble. Conditional on the true class, the log-odds of a positive label is
decomposed into class-specific model and prompt effects plus three item-level random
effects: an item-ambiguity effect shared by every configuration, an item$\times$model
effect and an item$\times$prompt effect; repeated decoding runs are conditionally
independent given these effects. We call this the crossed random-effects latent class
model (CRE-LCM). It nests the Dawid--Skene model and the one-factor random-effects model
of \citet{Qu1996} as special cases, and it returns, in addition to posterior
class probabilities, a decomposition of annotation noise into item, model, prompt and run
components on a common scale.

Our contributions are threefold. First, we show by simulation that under crossed
dependence of realistic size the Dawid--Skene estimator overstates specificity and
inflates prevalence, and that the bias grows monotonically with the strength of
dependence. Second, we analyse the identifiability of the CRE-LCM. The crossed
item$\times$model and item$\times$prompt effects are identified by the replication
structure of the design, but the shared item-ambiguity effect competes with the latent
class itself and produces a second, merged-class posterior mode. We compare five remedies
and show that the mode is avoided by initialising the sampler at the Dawid--Skene
solution, and removed altogether by anchoring a small random subset of items with known
labels. Third, we quantify, over a grid of dependence levels, sample sizes and numbers of
prompts, how CRE-LCM compares with majority vote and Dawid--Skene in prevalence bias,
accuracy and calibration of posterior labels, interval coverage and recovery of the
variance shares, including a misspecified data-generating process with heavy-tailed
random effects.

Section~\ref{sec:background} reviews the three literatures the paper connects.
Section~\ref{sec:model} defines the model, its interpretation and its estimation, and
presents the identifiability analysis. Section~\ref{sec:simulation} reports the
simulation study. Section~\ref{sec:application} outlines the application and
Section~\ref{sec:discussion} discusses design implications and limitations.
